\documentclass[12pt,a4paper]{article}
\usepackage{mathwizards-ingles}
\title{%
  \vspace{-1cm}
  {\Huge\bfseries\textcolor{mes2}{Mes 2 · Semana 5}}\\[0.3cm]
  {\Large Past Simple, Lectura Inicial y Conectores Básicos}\\[0.2cm]
  {\large\color{grissuave} Nivel A1 · Introducción a la Lectura con Graded Readers}\\[0.5cm]
  \rule{\textwidth}{1.5pt}
}
\author{}
\date{}

\begin{document}
\maketitle
\thispagestyle{empty}

\begin{cajanivel}
\textbf{Objetivo:} El estudiante comienza a leer textos simples (\textit{Graded Readers} nivel A1) en los Círculos de Lectura Virtuales. Se introduce el \textbf{Past Simple} como el segundo tiempo verbal clave, junto con \textbf{conectores básicos} para construir oraciones más largas. Esta es la semana donde se introduce formalmente el texto impreso.
\end{cajanivel}

% ═══════════════════════════════════════════════════════════════════════════════
\section{Gramática Emergente: Past Simple — Verbo \textit{To Be}}
% ═══════════════════════════════════════════════════════════════════════════════

\begin{cajagrammar}[Past Simple de \textit{to be}: was / were]
\begin{tabularx}{\textwidth}{@{} l l X @{}}
\toprule
\textbf{Sujeto} & \textbf{Forma} & \textbf{Ejemplo} \\
\midrule
I / He / She / It & \eng{was} & \textit{I \textbf{was} at home yesterday.} \\
You / We / They & \eng{were} & \textit{They \textbf{were} happy.} \\
\midrule
\multicolumn{3}{@{}l}{\textbf{Negativo:}} \\
I / He / She / It & \eng{wasn't} & \textit{She \textbf{wasn't} at school.} \\
You / We / They & \eng{weren't} & \textit{We \textbf{weren't} tired.} \\
\midrule
\multicolumn{3}{@{}l}{\textbf{Pregunta:}} \\
& \eng{Was} I / he / she / it...? & \textit{\textbf{Was} he at the party?} \\
& \eng{Were} you / we / they...? & \textit{\textbf{Were} you happy?} \\
\bottomrule
\end{tabularx}
\end{cajagrammar}

% ═══════════════════════════════════════════════════════════════════════════════
\section{Gramática Emergente: Past Simple — Verbos Regulares}
% ═══════════════════════════════════════════════════════════════════════════════

\begin{cajagrammar}[Past Simple Regular: verbo + \textbf{-ed}]
Se usa para hablar de \textbf{acciones completadas en el pasado}.

\medskip
\begin{tabularx}{\textwidth}{@{} l l l @{}}
\toprule
\textbf{Regla} & \textbf{Base} & \textbf{Pasado} \\
\midrule
General: + ed & \textit{walk} & \textit{walk\textbf{ed}} \\
Termina en -e: + d & \textit{live} & \textit{live\textbf{d}} \\
Termina en consonante + y: -y $\rightarrow$ -ied & \textit{study} & \textit{stud\textbf{ied}} \\
CVC corta (1 sílaba): doblar + ed & \textit{stop} & \textit{stopp\textbf{ed}} \\
\bottomrule
\end{tabularx}

\medskip
\begin{cajaejemplo}[Pronunciación del -ed]
El sufijo \textbf{-ed} tiene tres pronunciaciones:
\begin{itemize}[nosep, leftmargin=1cm]
  \item \textbf{/t/} después de sonidos sordos: \textit{walked, washed, cooked}
  \item \textbf{/d/} después de sonidos sonoros: \textit{lived, played, opened}
  \item \textbf{/ɪd/} después de /t/ o /d/: \textit{wanted, needed, started}
\end{itemize}
\end{cajaejemplo}

\medskip
\textbf{Negativo:} Sujeto + \eng{didn't} + verbo base\\
\textit{I \textbf{didn't watch} TV last night.} \quad ($\times$ I didn't watch\textbf{ed})

\medskip
\textbf{Pregunta:} \eng{Did} + sujeto + verbo base + ?\\
\textit{\textbf{Did} you \textbf{study} yesterday?} \quad ($\times$ Did you studi\textbf{ed})
\end{cajagrammar}

% ═══════════════════════════════════════════════════════════════════════════════
\section{Gramática Emergente: Verbos Irregulares de Alta Frecuencia}
% ═══════════════════════════════════════════════════════════════════════════════

\begin{cajagrammar}[Top 20 Verbos Irregulares — Memorizar como Chunks]
Estos son los verbos irregulares que aparecen con mayor frecuencia en los \textit{Graded Readers} A1. No siguen la regla del ``-ed''.

\medskip
\begin{multicols}{2}
\begin{tabularx}{\linewidth}{@{} l l l @{}}
\toprule
\textbf{Base} & \textbf{Past} & \textbf{Significado} \\
\midrule
go & \eng{went} & ir \\
have & \eng{had} & tener \\
get & \eng{got} & obtener \\
see & \eng{saw} & ver \\
come & \eng{came} & venir \\
make & \eng{made} & hacer \\
take & \eng{took} & tomar \\
give & \eng{gave} & dar \\
say & \eng{said} & decir \\
tell & \eng{told} & contar \\
\bottomrule
\end{tabularx}

\columnbreak

\begin{tabularx}{\linewidth}{@{} l l l @{}}
\toprule
\textbf{Base} & \textbf{Past} & \textbf{Significado} \\
\midrule
think & \eng{thought} & pensar \\
know & \eng{knew} & saber \\
find & \eng{found} & encontrar \\
eat & \eng{ate} & comer \\
drink & \eng{drank} & beber \\
run & \eng{ran} & correr \\
write & \eng{wrote} & escribir \\
read & \eng{read} & leer \\
buy & \eng{bought} & comprar \\
bring & \eng{brought} & traer \\
\bottomrule
\end{tabularx}
\end{multicols}
\end{cajagrammar}

% ═══════════════════════════════════════════════════════════════════════════════
\section{Chunks: Expresiones de Tiempo Pasado}
% ═══════════════════════════════════════════════════════════════════════════════

\begin{cajachunk}[Past Time Expressions — Marcadores Temporales del Pasado]
Estas expresiones actúan como ``señales'' de que la oración está en pasado:

\medskip
\begin{multicols}{2}
\begin{tabularx}{\linewidth}{@{} X l @{}}
\eng{yesterday} & \esp{ayer} \\
\eng{last night} & \esp{anoche} \\
\eng{last week} & \esp{la semana pasada} \\
\eng{last month} & \esp{el mes pasado} \\
\eng{last year} & \esp{el año pasado} \\
\end{tabularx}

\columnbreak

\begin{tabularx}{\linewidth}{@{} X l @{}}
\eng{two days ago} & \esp{hace dos días} \\
\eng{a week ago} & \esp{hace una semana} \\
\eng{this morning} & \esp{esta mañana} \\
\eng{on Monday} & \esp{el lunes} \\
\eng{in 2020} & \esp{en 2020} \\
\end{tabularx}
\end{multicols}
\end{cajachunk}

% ═══════════════════════════════════════════════════════════════════════════════
\section{Chunks: Conectores Básicos}
% ═══════════════════════════════════════════════════════════════════════════════

\begin{cajachunk}[Basic Connectors — Para Construir Oraciones Más Largas]
Estos cuatro conectores son esenciales para avanzar más allá de las oraciones simples:

\medskip
\begin{tabularx}{\textwidth}{@{} l l X @{}}
\toprule
\textbf{Conector} & \textbf{Significado} & \textbf{Ejemplo} \\
\midrule
\eng{and} & y (adición) & \textit{I went to the store \textbf{and} bought some food.} \\[4pt]
\eng{but} & pero (contraste) & \textit{I was tired, \textbf{but} I didn't go to bed.} \\[4pt]
\eng{because} & porque (causa) & \textit{She stayed home \textbf{because} she was sick.} \\[4pt]
\eng{so} & así que (resultado) & \textit{It was raining, \textbf{so} I took an umbrella.} \\
\bottomrule
\end{tabularx}

\medskip
\begin{cajaejemplo}[Patrón Mental]
\textbf{and} = ``y además''\\
\textbf{but} = ``sin embargo''\\
\textbf{because} = ``la razón es que...''\\
\textbf{so} = ``por lo tanto...''
\end{cajaejemplo}
\end{cajachunk}

% ═══════════════════════════════════════════════════════════════════════════════
\section{Chunks: Frases de Lectura (\textit{Graded Reader Patterns})}
% ═══════════════════════════════════════════════════════════════════════════════

\begin{cajachunk}[Patrones Narrativos que Aparecen en Graded Readers A1]
Al leer historias simples, estos patrones aparecen constantemente:

\medskip
\begin{tabularx}{\textwidth}{@{} X l @{}}
\eng{Once upon a time...} & \esp{Érase una vez...} \\
\eng{One day, ...} & \esp{Un día, ...} \\
\eng{A long time ago, ...} & \esp{Hace mucho tiempo, ...} \\
\eng{Then, ...} / \eng{After that, ...} & \esp{Luego, ... / Después de eso, ...} \\
\eng{Suddenly, ...} & \esp{De repente, ...} \\
\eng{Finally, ...} & \esp{Finalmente, ...} \\
\eng{The next day, ...} & \esp{Al día siguiente, ...} \\
\eng{At the end, ...} & \esp{Al final, ...} \\
\eng{... and they lived happily ever after.} & \esp{... y vivieron felices para siempre.} \\
\end{tabularx}
\end{cajachunk}

% ═══════════════════════════════════════════════════════════════════════════════
\section{Oraciones Listas para Minar}
% ═══════════════════════════════════════════════════════════════════════════════

\begin{cajamineria}[Oraciones \textit{i+1} — Semana 5]

\begin{enumerate}[leftmargin=1.2cm]
  \item \textit{I \underline{went} to the supermarket yesterday.}\\
    {\small\textbf{Objetivo:} el pasado irregular ``went'' (go $\rightarrow$ went).}

  \item \textit{She \underline{had} a big breakfast this morning.}\\
    {\small\textbf{Objetivo:} el pasado irregular ``had'' (have $\rightarrow$ had).}

  \item \textit{We watched a movie, \underline{but} it was boring.}\\
    {\small\textbf{Objetivo:} el conector ``but'' (pero) para contrastar ideas.}

  \item \textit{He didn't come to class \underline{because} he was sick.}\\
    {\small\textbf{Objetivo:} el conector ``because'' para dar razones.}

  \item \textit{I \underline{saw} a beautiful bird in the park.}\\
    {\small\textbf{Objetivo:} el pasado irregular ``saw'' (see $\rightarrow$ saw).}

  \item \textit{She cook\underline{ed} dinner and we ate together.}\\
    {\small\textbf{Objetivo:} reconocer la terminación regular ``-ed'' (cooked).}

  \item \textit{\underline{Once upon a time}, there was a little girl.}\\
    {\small\textbf{Objetivo:} la fórmula narrativa ``Once upon a time''.}

  \item \textit{Did you \underline{buy} anything at the store?}\\
    {\small\textbf{Objetivo:} estructura interrogativa con ``Did'' + verbo base.}

  \item \textit{It was raining, \underline{so} I stayed home.}\\
    {\small\textbf{Objetivo:} el conector ``so'' (así que) para consecuencias.}

  \item \textit{He \underline{told} me a funny story last night.}\\
    {\small\textbf{Objetivo:} el pasado irregular ``told'' (tell $\rightarrow$ told).}
\end{enumerate}
\end{cajamineria}

\vfill
\begin{center}
\rule{0.5\textwidth}{0.5pt}\\[0.3cm]
{\small\color{grissuave} Curso de Inmersión en Inglés · Mes 2 · Semana 5 · Nivel A1}
\end{center}

\end{document}
